\documentclass[11pt]{article}
\usepackage[T1]{fontenc}
\usepackage{graphicx}
\usepackage{amsmath,amssymb}
\usepackage{booktabs}
\usepackage{hyperref}
\usepackage{ragged2e}
\usepackage{microtype}
\usepackage{authblk}
\usepackage[margin=1in]{geometry}
\usepackage{enumitem}
\usepackage{array}
\justifying
\hypersetup{colorlinks=true,citecolor=blue,urlcolor=blue,linkcolor=blue}

\title{From Projected Subspaces to Physical Circuits:\\
Spin-Faithful ADAPT-VQE for an Open-Shell Copper Center}

\author[1]{Lucia Mal\'{\i}\v{c}kov\'a\thanks{Corresponding author: \href{mailto:lucy@thesafehaven.ai}{lucy@thesafehaven.ai}}}
\author[2,3]{Petr Klenovsk\'y\thanks{E-mail: \href{mailto:klenovsky@physics.muni.cz}{klenovsky@physics.muni.cz}}}
\affil[1]{Modelos Inteligencia Artificial, Cl. Tajinaste, San Miguel de Abona, Santa Cruz de Tenerife, Spain}
\affil[2]{Department of Condensed Matter Physics, Faculty of Science, Masaryk University, Kotl\'a\v{r}sk\'a 267/2, 61137 Brno, Czech Republic}
\affil[3]{Czech Metrology Institute, Okru\v{z}n\'i~31, 63800 Brno, Czech Republic}
\date{}

\begin{document}
\maketitle

\begin{abstract}
For open-shell quantum chemistry, an accurate variational energy is not sufficient evidence that the implemented quantum circuit prepares the intended physical state.
We demonstrate this distinction for an idealized type-1 (blue-copper) Cu(II) center using two exactly verifiable active-space Hamiltonians, CAS(15e,9o) and CAS(21e,12o), mapped to 18 and 24 qubits.
Both contain a lower quartet in addition to the target Cu(II)-motivated doublet, making explicit total-spin control essential.
An exact-doublet projected ADAPT-VQE calculation reaches a 1.572 m$E_h$ error with 123 operators, but applying the same optimized parameters to the corresponding unprojected fermionic generators gives only 0.074 fidelity with the projected state, $\langle S^2\rangle=2.382$, and a 464.7 m$E_h$ error.
We therefore replace the projected construction by generalized physical spin-preserving rank-1/rank-2 operator pools and use repeated ADAPT selection, global reoptimization, and pruning.
The resulting fermionic ans\"atze reach 1.493 m$E_h$ with 50 operator applications on 18 qubits and 1.181 m$E_h$ with 61 applications on 24 qubits, while maintaining $\langle S^2\rangle=0.75$ to numerical precision; their exact-state fidelities are 0.9826 and 0.9907.
Circuit synthesis exposes a second accuracy--resource tradeoff.
For the 24-qubit state, second-order Suzuki synthesis with one repetition reproduces the fermionic ansatz within 0.091 m$E_h$ and 0.99969 fidelity, yielding a total doublet error of 1.272 m$E_h$, but requires 22,574 CX gates at abstract-gate-set depth 26,409.
For the 18-qubit state, four Suzuki repetitions reduce the synthesis error to 0.135 m$E_h$ and raise the fidelity to 0.99948, at the cost of 61,482 CX gates and depth 71,636; the resulting total error is 1.628 m$E_h$.
Finally, ideal greedy qubit-wise-commuting measurement analysis requires approximately $1.43\times10^8$ and $5.46\times10^8$ optimally allocated shots, respectively, for a 1.6 m$E_h$ statistical uncertainty.
The benchmark therefore separates four validation layers---spin targeting, physical fermionic ansatz construction, qubit-circuit synthesis, and measurement cost---that should not be conflated when assessing VQE accuracy or hardware readiness.
\end{abstract}

\section{Introduction}
Variational quantum eigensolvers (VQEs) combine parameterized quantum circuits with classical optimization and remain a central approach to quantum electronic-structure simulation \cite{peruzzo2014vqe,mcclean2016theory,mcardle2020review}.
Their physical meaning, however, depends on several logically distinct layers: the molecular Hamiltonian, the target symmetry sector, the variational manifold, the physical fermionic unitary sequence, the synthesized qubit circuit, and the measurements used to estimate observables.
A favorable energy at one layer need not validate the next.

ADAPT-VQE grows an ansatz one operator at a time according to energy derivatives \cite{grimsley2019adapt}.
Generalized spin-complemented excitations and repeated operator selection already appear in the original method; later work has explored qubit pools \cite{tang2021qubitadapt}, symmetry breaking \cite{bertels2022symmetry}, compact circuit constructions \cite{magoulas2023cnot}, and pruning \cite{vaquero2025pruned}.
The present work therefore does \emph{not} claim generalized excitations, operator reuse, or pruning as new algorithms.
Instead, we address a validation problem that becomes acute for open-shell systems: a state optimized after exact projection into a target spin subspace need not be represented by the bare fermionic circuit whose resources are subsequently quoted.

Type-1, or blue-copper, sites provide a chemically motivated benchmark.
Their electronic structure is characterized by a strongly covalent Cu--thiolate interaction and an oxidized Cu(II) center with nominal $d^9$ doublet character; such sites mediate electron transfer in blue-copper proteins and multicopper oxidases \cite{solomon2004bluecopper,solomon2006thiolate,solomon2014copper}.
The model used here is deliberately small and idealized.
It is valuable precisely because exact spin-resolved classical references remain available, allowing every stage of the VQE pipeline to be falsified rather than assumed.

Our starting diagnostic uses the exact target-doublet subspace.
If $U$ spans that subspace, a projected generator $A_d=U^\dagger A U$ defines a valid reduced-space emulator evolution.
It does not, in general, imply that $e^{\theta A}$ restricted to the same state is equivalent to $e^{\theta A_d}$.
We quantify this mismatch, replace it with explicitly spin-preserving physical generators, scale the construction from 18 to 24 qubits, and then test the actual Jordan--Wigner/Suzuki circuits and their measurement burden.

Throughout, the commonly used 1.6 m$E_h$ threshold is employed only as an \emph{algorithmic accuracy benchmark for a specified active-space Hamiltonian}.
It is not a claim of 1 kcal/mol predictive accuracy for the real protein, for laccase, or for the underlying STO-3G model.

\section{Computational methodology}

\subsection{Copper-center Hamiltonians}
The benchmark uses an idealized oxidized type-1 Cu(II) first-coordination-shell model with two nitrogen donors, a short cysteine-model sulfur donor, and a longer axial sulfur donor.
The electronic-structure workflow was generated with PySCF \cite{sun2018pyscf}; the primary active space was selected with the atomic-valence active-space (AVAS) procedure \cite{sayfutyarova2017avas}, targeting the Cu $3d$ and cysteine-S $3p$ manifolds.

The primary Hamiltonian is CAS(15e,9o), yielding 18 spin orbitals.
An expanded CAS(21e,12o) Hamiltonian from the same model provides a 24-qubit algorithmic scaling benchmark.
Because the two active spaces were not generated as a formally balanced state-specific convergence sequence for both multiplicities, their comparison is used to test the robustness of the quantum algorithm, not to claim converged spin-state energetics.
Both Hamiltonians are mapped with the Jordan--Wigner transformation; alternative mappings such as Bravyi--Kitaev may reduce Pauli weights \cite{seeley2012bk}, but are not considered here.

\subsection{Exact spin-resolved references}
For each Hamiltonian, the target determinant sector has $M_S=1/2$ and fixed particle number.
These quantum numbers do not fix total spin.
Let $\hat S_+$ be the many-electron spin-raising operator.
Because the target has $M_S=S=1/2$, the pure-doublet highest-weight subspace is $\mathcal H_D=\ker \hat S_+$.
With $U$ an orthonormal basis of this null space, $H_D=U^\dagger H_{M_S}U$ provides an exact doublet oracle.

The target doublet is selected because the benchmark is intended to represent an oxidized Cu(II) center.
The active-space Hamiltonians themselves place a quartet below that doublet.
We therefore call it the \emph{non-target quartet}, rather than an unphysical eigenstate.
The calculated ordering is not interpreted as a converged prediction of the real blue-copper spin ordering.

\subsection{Projected-subspace equivalence test}
The diagnostic projected ADAPT calculation is performed entirely in the exact doublet basis using $A_{d,k}=U^\dagger A_k U$.
To test whether this reduced-space state corresponds to the physical bare sequence, we reconstruct
\begin{align}
|\Psi_{\rm proj}\rangle &= \prod_k e^{\theta_k A_{d,k}}|\Phi_{\rm HF}\rangle_D,\\
|\Psi_{\rm bare}\rangle &= \prod_k e^{\theta_k A_k}|\Phi_{\rm HF}\rangle
\end{align}
from the same 123 selected operators and optimized parameters.
We compare energy, $\langle S^2\rangle$, doublet weight, projected/bare fidelity, and generator leakage.

\subsection{Physical spin-preserving ADAPT}
All subsequent adaptive calculations use physical fermionic generators that preserve the target doublet directly.
Generalized rank-1 and rank-2 $M_S$-conserving determinant excitations are grouped by their spatial remove/add pattern.
Within each group, spin-preserving linear combinations are obtained from the numerical null space of the doublet-to-complement coupling.
A generator $B$ is retained only if its doublet leakage and normalized $S^2$ commutator defect are below numerical tolerance.
For the final 18-qubit pool, the maximum measured leakage ratio is $3.14\times10^{-16}$ and the maximum normalized $S^2$ commutator defect is $1.11\times10^{-15}$.

The final generalized pools contain 1,080 generators at 18 qubits and 3,762 at 24 qubits.
ADAPT selection uses the magnitude of the energy derivative
\begin{equation}
g_k=2\,{\rm Re}\langle\Psi|H B_k|\Psi\rangle.
\end{equation}
The selected generator is appended without draining the pool, allowing repetitions.
Parameters are reoptimized with L-BFGS-B, including periodic global reoptimization.
High-accuracy sequences are then pruned by removing small-amplitude applications and reoptimizing the remainder.
This pruning is used as a practical compaction step and is not presented as a new pruning algorithm \cite{vaquero2025pruned}.

\subsection{Qubit mapping and circuit synthesis}
The physical fermionic generators are mapped using Qiskit Nature 0.7.2 and Jordan--Wigner, and the resulting Hermitian Pauli operators are exponentiated with Qiskit 2.2.3 \texttt{PauliEvolutionGate} using a second-order Suzuki product formula \cite{qiskit2024,qiskitnature2023}.
Qiskit Aer 0.17.2 is used for statevector validation.
Circuits are transpiled at optimization level 1 to the abstract basis $\{R_z,\sqrt X,X,{\rm CX}\}$.
The quoted depths and CX counts are therefore \emph{not} device-mapped Euro-Q-Exa resource estimates.

Circuit validation includes full-state fidelity to the exact fermionic sequence, target-$M_S$ survival, conditioned $\langle S^2\rangle$, and synthesis energy error.
A strict circuit-equivalence gate requires fidelity $\ge 0.999$, $|\delta E_{\rm synth}|\le0.1$ m$E_h$, $|\langle S^2\rangle-0.75|\le0.001$, and target-$M_S$ survival $\ge0.999$.

\subsection{Measurement and wavefunction diagnostics}
The qubit Hamiltonians are grouped greedily into qubit-wise-commuting (QWC) sets.
For the chosen grouping, exact statevector variances are used to estimate the shot count under variance-optimal allocation,
\begin{equation}
N_{\rm shot}^{\rm opt}=
\frac{\left(\sum_g\sqrt{{\rm Var}_g}\right)^2}{\sigma_E^2}.
\end{equation}
This is the minimum for that fixed grouping and ideal state preparation; it is not a global lower bound over all possible measurement strategies and excludes device noise and calibration overhead \cite{verteletskyi2020measurement}.

Wavefunction quality is assessed independently using exact-state fidelity, the residual norm $r=\|(H-E)|\Psi\rangle\|$, one-particle reduced density matrices, and natural occupations.
The compact ans\"atze are also reoptimized from perturbed initial parameters to assess local robustness.

\section{Results}

\subsection{Exact target-doublet references}
Table~\ref{tab:references} summarizes the exact spin-resolved references.
Both active spaces place the non-target quartet below the target doublet by about 6 m$E_h$.
The similar values show that the lower quartet is not unique to the nine-orbital Hamiltonian, but they should not be read as chemically converged spin gaps.

\begin{table}[htbp]
\centering
\small
\caption{Exact spin-resolved references. The target is the lowest doublet, not the global active-space eigenstate.}
\label{tab:references}
\begin{tabular}{lcc}
\toprule
& CAS(15e,9o) & CAS(21e,12o)\\
\midrule
Qubits & 18 & 24\\
$M_S=1/2$ determinant dimension & 324 & 792\\
Exact doublet dimension & 240 & 572\\
$E_D$ (Hartree) & $-2518.989067370$ & $-2518.989131438$\\
$E_Q$ (Hartree) & $-2518.995009103$ & $-2518.995161622$\\
$E_D-E_Q$ (m$E_h$) & 5.941733 & 6.030184\\
HF determinant energy (Hartree) & $-2518.540295701$ & $-2518.754566933$\\
\bottomrule
\end{tabular}
\end{table}

\subsection{Projected accuracy does not imply physical-circuit accuracy}
The projected doublet ADAPT calculation reaches $E_{\rm proj}=-2518.987495837~E_h$ and $\Delta E_{\rm proj}=1.571533$ m$E_h$, with exactly $\langle S^2\rangle=0.75$.
The corresponding bare physical sequence, using the same selected operators and angles, instead gives $E_{\rm bare}=-2518.524355925~E_h$, or 464.711 m$E_h$ above the target doublet, with $\langle S^2\rangle=2.381596$.
Its doublet weight is 0.4561 and its fidelity with the projected state is only 0.07416.
Moreover, 116 of the 123 selected bare generators exhibit doublet leakage above $10^{-10}$.
The discrepancy is therefore structural, not a small numerical or transpilation error.

\begin{figure}[htbp]
\centering
\includegraphics[width=0.98\linewidth]{fig1_method_ablation.pdf}
\caption{Method ablation. Exact spin projection produces an accurate reduced-space state, but the same bare generator sequence fails. Physical spin adaptation cures the symmetry problem; generalized excitations and repeated selection are then required to recover sufficient expressivity.}
\label{fig:ablation}
\end{figure}

\subsection{Spin correctness alone is not enough}
The first physical spin-adapted pool is restricted to Hartree--Fock occupied-to-virtual excitations.
With S+D generators, all 35 operators are exhausted at 153.629 m$E_h$ error.
Adding triples gives 147 physical generators but reaches only 141.392 m$E_h$.
Both states remain exactly spin pure.

Allowing repeated selection from the same restricted S+D+T pool improves the 250-application state to 13.166 m$E_h$ error, fidelity 0.8783, and residual 0.0412 $E_h$, while preserving $\langle S^2\rangle=0.75$.
This is a substantial improvement, but it still misses the 1.6 m$E_h$ model-space benchmark by almost an order of magnitude.
The result isolates the role of pool expressivity: repeats help, but a generalized physical pool is still required.

\subsection{Generalized physical pools recover compact accurate states}
At 18 qubits, the 1,080-generator generalized physical pool yields a compact state with 50 operator applications, 39 of them unique.
Its energy error is 1.492674 m$E_h$, with exact-state fidelity 0.982614, residual 0.048011 $E_h$, and unit doublet weight to numerical precision.

At 24 qubits, the 3,762-generator pool yields a 61-application state using 41 unique generators.
Its error is 1.181293 m$E_h$, with fidelity 0.990681 and residual 0.058536 $E_h$.
Again $\langle S^2\rangle=0.75$ and the doublet weight is unity to numerical precision.

\begin{table}[htbp]
\centering
\small
\caption{Ablation of spin treatment, operator pool, and reuse. Errors refer to the target doublet.}
\label{tab:ablation}
\begin{tabular}{p{5.1cm}rrrr}
\toprule
Method & Ops & $\Delta E$ (m$E_h$) & $\langle S^2\rangle$ & Fidelity\\
\midrule
Projected-subspace ADAPT, 18q & 123 & 1.572 & 0.750 & ---\\
Same angles, bare generators, 18q & 123 & 464.711 & 2.382 & 0.074$^\dagger$\\
Restricted physical S+D, 18q & 35 & 153.629 & 0.750 & 0.2866\\
Restricted physical S+D+T, 18q & 147 & 141.392 & 0.750 & 0.2672\\
Restricted S+D+T + repeats, 18q & 250 & 13.166 & 0.750 & 0.8783\\
Generalized physical + repeats, 18q & 50 & 1.493 & 0.750 & 0.9826\\
Generalized physical + repeats, 24q & 61 & 1.181 & 0.750 & 0.9907\\
\bottomrule
\end{tabular}

\raggedright\footnotesize
$^\dagger$Fidelity to the projected-subspace state. Other fidelities are to the exact target doublet.
\end{table}

The compact solutions are locally robust to parameter perturbations.
Seven 18-qubit reoptimizations finish in the 1.49266--1.49306 m$E_h$ range, while five 24-qubit reoptimizations finish in the 1.18070--1.18125 m$E_h$ range (Fig.~\ref{fig:multistart}).
Some individual optimizer runs reach their iteration cap, so this test should be interpreted as energetic local robustness rather than a proof of global optimization.

\begin{figure}[htbp]
\centering
\includegraphics[width=0.83\linewidth]{fig5_multistart.pdf}
\caption{Multi-start robustness of the compact generalized physical ans\"atze. Perturbed parameter vectors return to narrow energy windows while preserving the target spin.}
\label{fig:multistart}
\end{figure}

\subsection{Circuit synthesis introduces a second approximation}
The physical fermionic ansatz and the mapped qubit circuit now use the same generators.
Product-formula synthesis nevertheless introduces an additional approximation.

For the 18-qubit 50-application state, one second-order Suzuki repetition gives a 44.381 m$E_h$ synthesis error and circuit-to-fermionic fidelity 0.8454.
Two repetitions reduce these values to 2.416 m$E_h$ and 0.99136.
Four repetitions reach 0.135 m$E_h$ synthesis error and 0.999484 fidelity.
The total error relative to the exact target doublet is then 1.628 m$E_h$, slightly above the 1.6 m$E_h$ model-space benchmark.
The four-repetition circuit does \emph{not} pass our strict 0.1 m$E_h$ synthesis criterion.

For the 24-qubit 61-application state, one Suzuki repetition already passes all strict circuit-equivalence gates:
the synthesis error is 0.0906 m$E_h$, fidelity is 0.999694, conditioned $\langle S^2\rangle=0.750624$, and target-$M_S$ survival exceeds 0.999999999.
Its total target-doublet error is 1.272 m$E_h$.
This favorable product-formula behavior is instance specific and is not evidence of favorable asymptotic scaling.

\begin{table}[htbp]
\centering
\small
\caption{Circuit synthesis results after transpilation to an abstract $\{R_z,\sqrt X,X,\mathrm{CX}\}$ basis. These are not device-mapped resources.}
\label{tab:circuit}
\begin{tabular}{lrrrrrr}
\toprule
Case & Suzuki reps & CX & Depth & $\delta E_{\rm synth}$ & $F_{\rm circ}$ & Total $\Delta E$\\
& & & & (m$E_h$) & & (m$E_h$)\\
\midrule
18q, 50 ops & 1 & 15,936 & 18,389 & 44.381 & 0.84541 & 45.873\\
18q, 50 ops & 2 & 31,118 & 36,138 & 2.416 & 0.99136 & 3.909\\
18q, 50 ops & 4 & 61,482 & 71,636 & 0.135 & 0.99948 & 1.628\\
24q, 61 ops & 1 & 22,574 & 26,409 & 0.091 & 0.99969 & 1.272\\
\bottomrule
\end{tabular}
\end{table}

\begin{figure}[htbp]
\centering
\includegraphics[width=0.83\linewidth]{fig2_synthesis_tradeoff_18q.pdf}
\caption{Second-order Suzuki synthesis of the 18-qubit 50-application ansatz. Increasing the repetition count rapidly improves state and energy fidelity, but approximately doubles the CX count each time.}
\label{fig:synthesis}
\end{figure}

The first five operators provide an implementation regression test:
with one Suzuki repetition, both 18- and 24-qubit circuits have circuit-to-fermionic fidelity above 0.999999998 and essentially exact spin purity.
Their energies remain hundreds of m$E_h$ above the target doublet, so these short prefixes validate implementation, not chemical accuracy.

For comparison, the older restricted 147-operator S+D+T physical sequence is itself synthesized accurately with one Suzuki repetition, but requires 218,502 CX gates at depth 261,514 while its fermionic energy error remains 141.392 m$E_h$.
This reinforces the need to optimize both variational expressivity and circuit resources.

\subsection{Measurement cost is already severe}
The 18-qubit Hamiltonian contains 9,316 Pauli terms and 1,909 greedy QWC groups.
The 24-qubit Hamiltonian contains 29,737 Pauli terms and 6,233 groups.
For the chosen groupings and ideal variance-optimal allocation, a 1.6 m$E_h$ statistical standard error requires approximately $1.43\times10^8$ shots at 18 qubits and $5.46\times10^8$ at 24 qubits.
At 0.5 m$E_h$, the corresponding estimates are $1.46\times10^9$ and $5.59\times10^9$ shots.

\begin{table}[htbp]
\centering
\small
\caption{Ideal shot estimates for the chosen greedy QWC grouping and optimized physical states. These are not global lower bounds over all measurement strategies and are not QPU runtime predictions.}
\label{tab:measurement}
\begin{tabular}{lrrrrr}
\toprule
Case & Pauli terms & QWC groups & \multicolumn{3}{c}{Optimal shots}\\
& & & 1.6 m$E_h$ & 1.0 m$E_h$ & 0.5 m$E_h$\\
\midrule
18 qubits & 9,316 & 1,909 & $1.43\times10^8$ & $3.65\times10^8$ & $1.46\times10^9$\\
24 qubits & 29,737 & 6,233 & $5.46\times10^8$ & $1.40\times10^9$ & $5.59\times10^9$\\
\bottomrule
\end{tabular}
\end{table}

\begin{figure}[htbp]
\centering
\includegraphics[width=0.81\linewidth]{fig3_measurement_shots.pdf}
\caption{Ideal shot burden for the chosen greedy QWC groupings. Hardware noise, state-preparation error, calibration, and alternative measurement strategies are not included.}
\label{fig:shots}
\end{figure}

A qubit count that fits on a device is therefore insufficient evidence of hardware feasibility.
At the m$E_h$ scale, these models are measurement-limited even before device noise is included.
A near-term QPU experiment should be framed primarily as circuit/symmetry/noise validation rather than as a modest-shot demonstration of model-space chemical accuracy.

\subsection{Wavefunction diagnostics beyond the energy}
The exact-state fidelities of the compact physical states are high but not unity, and their residual norms remain 0.0480 and 0.0585 $E_h$.
The one-particle diagnostics nevertheless remain close to the exact target doublets.
For 18 qubits, the spatial 1-RDM Frobenius error is 0.1860, the maximum natural-occupation error is 0.002109, and the natural-occupation-vector $L_2$ error is 0.002983.
For 24 qubits the corresponding values are 0.1250, 0.001605, and 0.002985.

\begin{figure}[htbp]
\centering
\includegraphics[width=0.81\linewidth]{fig4_natural_occupation_errors.pdf}
\caption{Absolute natural-orbital occupation errors relative to the exact target doublet. Natural occupations are well reproduced even though the full residual norms are not negligible.}
\label{fig:natural}
\end{figure}

\section{Discussion}

\subsection{What is and is not new}
Generalized spin-complemented ADAPT pools, repeated operator selection, symmetry-aware adaptive VQE, and pruning all have prior literature \cite{grimsley2019adapt,bertels2022symmetry,vaquero2025pruned}.
The contribution here is therefore not a claim to those algorithmic ingredients individually.
The methodological result is the explicit end-to-end separation of four layers that can otherwise be conflated:
\begin{enumerate}[label=(\roman*)]
\item the target spin sector of the electronic Hamiltonian;
\item the variational state optimized in a reduced or full emulator space;
\item the physical fermionic generator sequence;
\item the synthesized qubit circuit and its measurement protocol.
\end{enumerate}
The projected counterexample shows that layer (ii) can be accurate while failing to define layer (iii).
The physical generalized construction forces layers (ii) and (iii) to coincide.
The Qiskit tests then quantify the approximation between layers (iii) and (iv).

\subsection{Spin symmetry is necessary but insufficient}
The restricted S+D and S+D+T calculations demonstrate that exact spin preservation alone does not guarantee a useful variational manifold.
Their energies remain more than 140 m$E_h$ above the target despite perfect doublet purity.
Repeated selection from the restricted pool recovers most of that missing flexibility, reducing the error to 13.166 m$E_h$, but generalized physical one- and two-body rotations are still needed to cross the 1.6 m$E_h$ model-space benchmark.

Conversely, the projected-subspace result demonstrates that low energy alone does not prove circuit validity.
For open-shell VQE studies, reporting energy and particle number without total-spin diagnostics and a physical-generator equivalence test can therefore be insufficient.

\subsection{Hardware realism}
Neither final long circuit is near-term hardware friendly.
The cleaner 24-qubit circuit still has 22,574 CX gates and depth 26,409 before device connectivity, native-gate restrictions, calibration, or error mitigation are considered.
The 18-qubit circuit requires still more gates when Suzuki error is pushed close to the fermionic-ansatz scale.
More efficient circuit decompositions, including CNOT-efficient excitation circuits \cite{magoulas2023cnot}, are a natural next direction.

Measurement is at least as important as state preparation.
The $10^8$--$10^9$ ideal shot estimates imply that a realistic hardware campaign should not be judged by whether a modest number of measurements reproduces the final energy to 1.6 m$E_h$.
More informative near-term observables include particle-number and $M_S$ survival, selected spin correlators, short-prefix state preparation, and error growth with circuit length.
No physical-QPU result is reported in the present work.

\subsection{Chemical interpretation and limitations}
The benchmark employs STO-3G-based active-space Hamiltonians and an idealized first-shell model.
The 1.6 m$E_h$ threshold is therefore an internal algorithmic benchmark, not an uncertainty estimate for a real enzyme.
Likewise, the lower quartet in both Hamiltonians should not be interpreted as evidence that the experimental oxidized T1 center has quartet ground-state character.
A balanced physical spin-state calculation would require state-specific orbital optimization for both multiplicities, systematic active-space and basis convergence, and a more realistic environment.

A preliminary Cu--S geometry extension was generated during development.
It failed the pre-specified CASSCF convergence and active-orbital continuity gates and is deliberately excluded from the scientific results.
The five same-Hamiltonian exact diagonalizations generated in that failed scan all retain a lower quartet, but because the underlying orbital-generation workflow did not pass validation, those points are not used as a geometry-dependent physical claim.

The 24-qubit CAS(21e,12o) calculation is similarly treated as an algorithmic scaling benchmark rather than as proof that the active space is converged chemically.

\section{Conclusions}
An exact-doublet projected ADAPT calculation can be numerically accurate without defining the physical circuit that is subsequently resource-counted.
For the present open-shell copper model, the projected state reaches 1.572 m$E_h$, whereas the same optimized parameters applied to bare generators yield a strongly spin-contaminated state with only 0.074 projected-state fidelity and a 464.7 m$E_h$ error.

A generalized physical spin-preserving construction removes this ambiguity.
Repeated adaptive selection, global reoptimization, and pruning produce 50- and 61-application fermionic ans\"atze at 18 and 24 qubits with errors of 1.493 and 1.181 m$E_h$, respectively, while preserving the target doublet exactly.
One-particle observables and multi-start tests support the internal consistency of these states.

The qubit-circuit layer remains costly.
At 24 qubits, one second-order Suzuki repetition passes the strict circuit-equivalence gates and gives a total 1.272 m$E_h$ error, but requires 22,574 CX gates.
At 18 qubits, four repetitions approach state equivalence at 61,482 CX gates while leaving a total error of 1.628 m$E_h$ and narrowly failing the chosen strict synthesis-energy criterion.
Ideal grouped measurements further require hundreds of millions of shots for m$E_h$ statistical precision.

The broader lesson is that open-shell VQE validation should separately demonstrate the intended spin sector, equivalence between optimized and physical fermionic generators, controlled qubit synthesis, and a realistic measurement budget.
Without all four checks, ``accurate VQE'' and ``hardware-ready circuit'' are not interchangeable statements.

\section*{Data and Code Availability}
The physical generalized-spin ADAPT implementation, compact 18- and 24-qubit ansatz checkpoints, FCIDUMP Hamiltonians, validation summaries, circuit-synthesis outputs, figure source data, and reproducibility scripts are openly available at
\url{https://github.com/lucia-malickova/New-VQE-validation}.
The superseded projected-ADAPT repository state is retained under the \texttt{legacy/} tree for provenance and as the explicitly tested failure case.
A tagged archival release corresponding exactly to the submitted manuscript will be created at submission.

\section*{Acknowledgements}
P.K. acknowledges funding from the European Innovation Council Pathfinder program under Grant Agreement No.~101185617 (QCEED), funding by the Institutional Subsidy for the Long-Term Conceptual Development of a Research Organization granted to the Czech Metrology Institute by the Ministry of Industry and Trade of the Czech Republic, and by the project \emph{Quantum Materials for Applications in Sustainable Technologies}, CZ.02.01.01/00/22\_008/0004572.
The authors acknowledge the EuroHPC access award to the Euro-Q-Exa infrastructure under project proposal No.~EHPC-QCP-2026Q01-017.
% Retain the Euro-Q-Exa sentence only if the submitted work actually used, or was materially supported by, that awarded access.

\bibliographystyle{unsrt}
\bibliography{references}
\end{document}